\par\medskip
\noindent
\underline{2. Partes fija y tórica de $T_\ell(A_K)$. Criterios de buena
reducción. Teorema de ortogonalidad para $\ell\ne p$.}

\medskip
\noindent
2.1. Como todo grupo algebraico conmutativo sobre $k$, $A_o$ posee un
mayor subtoro, cuya formación conmuta con toda extensión del cuerpo de
base [SGA~3 XVII~7.2.1, XII~1.12]:
\begin{equation}\tag{2.1.1}
 T_o\subset A_o.
\end{equation}
Cuando $k$ es perfecto, $T_o$ se caracteriza por la condición de que
$A_o^o/T_o$ sea una extensión de un esquema abeliano $B_o$ sobre $k$
por un grupo unipotente liso conexo [25]. Este último posee
necesariamente una serie de composición cuyos cocientes son todos
isomorfos a $\mathbb G_{ak}$ (SGA~3 XVII~4.1.3), y se tiene una
sucesión exacta
\begin{equation}\tag{2.1.2}
 0\longrightarrow U_o\longrightarrow A_o^o/T_o
 \longrightarrow B_o\longrightarrow0.
\end{equation}
En cualquier caso, sea o no $k$ perfecto, \underline{si} $A_o^o$ es
una extensión de un esquema abeliano $B_o$ por un esquema afín liso
conexo $L_o$,
\begin{equation}\tag{2.1.3}
 0\longrightarrow L_o\longrightarrow A_o^o
 \longrightarrow B_o\longrightarrow0,
\end{equation}
se sabe que esta estructura de extensión es única [25]. Además,
$T_o\subset L_o$, y $T_o$ es el mayor subtoro de $L_o$. Se tiene así
la sucesión exacta (2.1.2), con
\begin{equation}\tag{2.1.4}
 U_o=L_o/T_o.
\end{equation}

\medskip
\noindent
2.1.5. Observemos que $A_o^o$ posee \underline{rango unipotente cero}
(es decir, tras pasar a una clausura algebraica de $k$, no posee ningún
subgrupo isomorfo a $\mathbb G_a$) si y solo si $A_o^o/T_o$ es un
esquema abeliano $B_o$. En tal caso, las hipótesis precedentes se
aplican con $U_o=0$.

Sea $\ell$ un primo distinto de la característica de $k$, de modo que
$A_o^o$ es $\ell$-divisible. La inclusión (2.1.1) da entonces una
sucesión exacta de proesquemas en grupos:
\begin{equation}\tag{2.1.6}
\begin{tikzcd}[row sep=large,column sep=large]
0 \arrow[r] & T_\ell(T_o) \arrow[r]
 & T_\ell(A_o^o) \arrow[r] \arrow[d,"\wr"]
 & T_\ell(A_o^o/T_o) \arrow[r] \arrow[d,"\wr"] & 0\\
 & & T_\ell(A_o) & T_\ell(A_o/T_o) & {,}
\end{tikzcd}
\end{equation}
cuyos términos pueden interpretarse también como haces $\ell$-ádicos
libres constantes torcidos. Cuando se dispone de una sucesión exacta
(2.1.2) ---por ejemplo, cuando $k$ es perfecto o cuando $A_o^o$ posee
rango unipotente cero---, se tiene
\begin{equation}\tag{2.1.7}
 T_\ell(A_o^o/T_o)\xrightarrow{\ \sim\ }T_\ell(B_o),
\end{equation}
de modo que la sucesión exacta (2.1.6) se convierte en
\begin{equation}\tag{2.1.8}
 0\longrightarrow T_\ell(T_o)\longrightarrow T_\ell(A_o^o)
 \longrightarrow T_\ell(B_o)\longrightarrow0.
\end{equation}
Pongamos
\begin{equation}\tag{2.1.9}
 n=\dim A_K=\dim A_o,\qquad
 \mu=\dim T_o,\qquad
 \lambda=\dim U_o,\qquad
 \alpha=\dim B_o.
\end{equation}
Entonces
\begin{equation}\tag{2.1.10}
 n=\mu+\lambda+\alpha,
\end{equation}
y
\begin{equation}\tag{2.1.11}
 \left\{
 \begin{aligned}
 \operatorname{rank}T_\ell(T_o)&=\mu,&
 \operatorname{rank}T_\ell(A_o^o/T_o)
 &=\operatorname{rank}T_\ell(B_o)=2\alpha,\\
 \operatorname{rank}T_\ell(A_o^o)&=\mu+2\alpha.
 \end{aligned}
 \right.
\end{equation}
La segunda relación es un conocido teorema de Weil [5, 15].

\medskip
\noindent
2.1.12. Los enteros $n,\mu,\lambda,\alpha$ están manifiestamente
definidos, tras pasar a una clausura algebraica de $k$, incluso sin
suponer la existencia de una sucesión exacta (2.1.2). Se denominan,
respectivamente, dimensión de $A_o$, \underline{rango reductivo de
$A_o$} (igual a la dimensión de $T_o$), \underline{rango unipotente de
$A_o$} y \underline{rango abeliano de $A_o$}. Siguen ligados por
(2.1.10) y (2.1.11), omitiendo de esta última el término
$\operatorname{rank}T_\ell(B_o)$.

\medskip
\noindent
2.1.13. Supongamos que $k$ posee característica $p>0$ y que $A_o^o$ es
$p$-divisible, o equivalentemente que posee rango unipotente cero
(cf. 2.1.5). Debe considerarse entonces también el progrupo
$T_p(A_o^o)$, que es un \underline{grupo de Barsotti--Tate} (o un
\underline{grupo $p$-divisible}, en la terminología de [33]). La
sucesión exacta (2.1.8) sigue siendo válida para $\ell=p$ y es entonces
una sucesión exacta de grupos de Barsotti--Tate. Aquí $\lambda=0$ y,
por tanto,
\begin{equation}\tag{2.1.13.1}
 n=\mu+\alpha.
\end{equation}
Las fórmulas (2.1.11) siguen siendo válidas asimismo, siempre que el
rango se interprete como el rango (o «altura») de los grupos de
Barsotti--Tate en cuestión.

\medskip
\noindent
2.1.14. Para ${A'}_o^o$ introducimos análogamente los invariantes
\[
 T'_o\subset {A'}_o^o,\qquad U'_o,\qquad B'_o,
\]
donde $T'_o$ es la parte tórica de ${A'}_o^o$, $B'_o$ su parte abeliana
(definida, en particular, cuando $k$ es perfecto o cuando ${A'}_o^o$
posee rango unipotente cero), y $U'_o$ su parte unipotente (definida,
por ejemplo, cuando $k$ es perfecto, o trivialmente cuando
${A'}_o^o$ posee rango unipotente cero). Existe una sucesión exacta
correspondiente de la forma (2.1.6), o (2.1.8), que no volveremos a
escribir. Podrían introducirse enteros $\mu',\lambda',\alpha'$ para
${A'}_o^o$ como en (2.1.9), pero veremos más adelante (2.2.7) que
coinciden con los de $A_o^o$; esto nos evita introducir para ellos una
notación nueva.
